\section{Ethical Considerations}
\label{sec:ethics}

Research on antisemitism detection raises important ethical questions that merit explicit discussion.

\paragraph{Dual use and censorship risk.}
Automated hate-speech detection systems carry inherent dual-use risk. A system designed to identify antisemitic content for moderation could be repurposed to suppress legitimate political speech, particularly criticism of Israeli government policies. We emphasize that our system is designed as a \emph{triage tool} to assist human content moderators, not as an autonomous decision-maker. The final determination of whether content violates platform policies should always involve human judgment, particularly for cases near the criticism--hate boundary. We recommend against using predicted probabilities below 0.8 for automated action, and instead directing such cases to human review queues.

\paragraph{Annotation framework and political context.}
Our dataset uses the International Holocaust Remembrance Alliance (IHRA) Working Definition of Antisemitism~\citep{ihra2016working} as its ground truth, following the dataset authors. The IHRA definition is widely adopted by governments and institutions, but it is also academically and politically contested: the Jerusalem Declaration on Antisemitism~\citep{jd2021jerusalem} offers an alternative framework that draws the criticism--hate boundary differently, particularly regarding anti-Zionist speech. We chose IHRA because it is the annotation framework of GoldStandard2024~\citep{jikeli2024goldstandard}; importantly, our methodological contribution (CCI v2, the multicalibration tension finding, and the structural Lemma~\ref{thm:ts-cfr}) is \emph{independent of the annotation framework} --- the method enforces counterfactual invariance with respect to whichever labels the annotators provide. Practitioners adopting our method under a different framework (e.g., Jerusalem Declaration) would inherit that framework's boundary judgments rather than IHRA's. We do not take a position in the IHRA debate but transparently acknowledge this dependency and its implications for model outputs.

\paragraph{Identity-term bias and counterfactual fairness.}
A core ethical concern is whether our models disproportionately flag benign content that mentions Jewish identity. If a model learns to associate ``Jews'' or ``Jewish'' with antisemitism, it will generate false positives on neutral or positive discussions of Jewish culture, religion, and community --- effectively penalising Jewish-topic speech. This is the precise concern CCI v2 is designed to address: by enforcing that the model produces the same prediction on a sentence and on its identity-token-swapped counterfactual, we discourage the model from using the identity term itself as a label-predictive feature. The FNED audit (\S\ref{sec:fairness}) measures false-negative-rate disparity across keyword groups; we report a Pareto trade-off explicitly --- CCI v2 (fixed T, JS, our pivot) reduces the counterfactual flip rate substantially at base scale, with only a mild FNED increase versus the strongest CLP baseline (Davani $\lambda=0.5$, $+2.9\%$, within one standard deviation) and an actual FNED \emph{improvement} versus same-architecture Model~C ($-7.0\%$ on base, $-1.5\%$ on DeBERTa-v3-large). The pre-registered CCI v2 (learned T, JS) cell, retained as an ablation comparator, has a much larger FNED regression ($+14.6\%$ vs.\ Davani); we report this transparently in \S\ref{sec:disc_fned}. Practitioners with FNED as a hard constraint should select the Davani $\lambda \in [0.5, 1.0]$ recipe; practitioners optimising for counterfactual invariance should use CCI v2.

\paragraph{Counterfactual swap pairs at training and evaluation time.}
For methodological transparency we note that counterfactual identity-token swaps are used at \emph{both} stages of our pipeline: at training time (as paired forward passes for the CCI v2 loss term, eq.~\ref{eq:cci-loss}) and at evaluation time (as a fixed 2{,}117-pair manifest over the test set for computing CFR). The swap engine is fully algorithmic --- a heuristic symmetry classifier rejects unsafe swaps involving slurs or geopolitical proxies --- and uses no human annotation; it produces transformations of the existing GoldStandard2024 data rather than newly collected data. We do not consider this dataset creation, but we disclose the dual training-time / evaluation-time use explicitly.

\paragraph{Dataset sensitivity and researcher welfare.}
The GoldStandard2024 dataset contains antisemitic content including slurs, conspiracy theories, and dehumanizing language that may be distressing to researchers and reviewers. We handle this data with care, reproduce offensive content only when necessary for scientific clarity, and apply content warnings in code documentation. Exposure to hateful content can have cumulative psychological effects on researchers and content moderators; institutions supporting work in this area should provide appropriate support.

\paragraph{Computing infrastructure and third-party inference.}
All fine-tuned-transformer training and the Mistral-7B-Instruct, Qwen-2.5-7B-Instruct, and Llama-3.1-8B-Instruct prompted inference are run on the EPFL Research Computing Platform (RCP) under our group's allocation. The Llama-3.3-70B-Versatile evaluation is performed via the Groq inference API (free tier) because the model is impractical to host locally; this means a subset of test-set tweets is sent to a third-party (US-based) inference provider. GoldStandard2024 tweets are public-domain content under CC BY 4.0, so no privacy violation results, but we disclose the data flow explicitly.

\paragraph{Reproducibility and open science.}
All code, configuration files, model architectures, random seeds, and the two-family Holm-Bonferroni-corrected paired-bootstrap significance suite (54 main-family tests + 6 within-architecture transfer tests) are publicly available. The GoldStandard2024 dataset is accessible under a CC BY 4.0 license from Zenodo~\citep{jikeli2024goldstandard}. We report all hyperparameters, preprocessing steps, and evaluation procedures in sufficient detail for independent reproduction. Model weights can be reproduced by re-running our training scripts with the provided configuration files; pre-trained checkpoints will be released on acceptance under the same licence as the dataset.

\paragraph{Scope of claims.}
We evaluate our models on a specific dataset (GoldStandard2024) with specific annotations (IHRA-based) in a specific language (English) from a specific platform (Twitter/X). We do not claim that our results generalize to all forms of antisemitism, all languages, all platforms, or all annotation frameworks. We additionally study only one identity target (antisemitism); we explicitly do not claim CCI v2 transfers to other identity-based hate-speech tasks without separate empirical validation, although the method is, in principle, identity-agnostic.
